\begin{adjustwidth}{-1cm}{-1cm}
%% Abstract (configurado para língua inglesa)
\begin{resumo}[Abstract]			% Título do Resumo (Abstract = Resumo em inglês)
\begin{otherlanguage*}{english}		% Língua do texto

The development of mobile robotic platforms in educational and research environments demands flexible and accessible architectures. This work proposes the adaptation of a commercial vacuum robot (Roomba® 614 iRobot) into an autonomous line-follower robotic platform. The system employs a Raspberry Pi 4 Model B, running Ubuntu 22.04, as the central processing unit. This choice eliminates the need for intermediary microcontrollers, allowing direct reading of a 5-channel line-follower module (BFD-1000) and the execution of the Proportional-Integral-Derivative (PID) control algorithm via Edge Computing. The entire control and communication software was developed modularly on the Robot Operating System 2 (ROS 2 Humble) ecosystem using C++ and Python. The integration of these technologies resulted in a simplified and robust hardware architecture, utilizing a USB-Serial TTL converter cable to mitigate communication noise. The results demonstrate the physical validation of the kinematic model and the platform's potential as an advanced educational resource for teaching kinematics, automation, and embedded systems.

\vspace{\onelineskip}
\noindent
\textbf{Keywords}: mobile robotics; line follower; ROS 2 Humble; Raspberry Pi 4; PID control.
\end{otherlanguage*}
\end{resumo}
\end{adjustwidth}